\subsection{The adaptive hybrid representation}
% ---- NOTES: The adaptive hybrid representation ----
% Takeaway: A hybrid is up to a budget of features selected from several base representations of the same molecules, chosen by a search for one dataset and one downstream model.
% Takeaway: Filtering out useless features matters as much as picking useful ones, because pooled irrelevant features act as noise.
% Paragraph 1: no single representation is best for every dataset, even for one model
% - Change: name the method "adaptive hybrid representation", "hybrid" for short. Drop "KIRBy (Knowledge Integration by Representation)" and "framework".
% - Change: "the noise robustness study above" should read "the noise robustness study in Chapter~\ref{chap:robustness}".
% - Example model: NGBoost, Gaussian noise, from that study's own results.
% - QM9 clean R² on held-out molecules (median of 10 replicates): PDV 0.865 is highest, MHG-GNN 0.852 is close behind, ECFP4 0.706 is lowest.
% - logD clean R² (all 5 folds for every representation): PDV 0.673 is highest, MHG-GNN 0.658 second.
% - hERG Ki clean R²: MHG-GNN 0.511 is highest (all 5 folds). PDV 0.456 is second, the median of the 4 folds whose clean R² was at least 0.3.
% - So NGBoost's best representation changes with the dataset: PDV on QM9 and logD, MHG-GNN on hERG Ki.
% - Caco-2 (log efflux ratio): Avalon 0.402 and ChemBERTa 0.404 are medians of only 2 of 5 folds. The other 3 folds fell below clean R² 0.3 and were excluded.
% - Over all 5 folds, Caco-2 reverses. Avalon's median clean R² is 0.287, ChemBERTa's 0.256, and PDV's 0.370.
% - Change: the Caco-2 claim, that NGBoost prefers Avalon, holds only after 3 of Avalon's 5 folds are dropped.
% - Change: "only strong with PDV on QM9" overstates it. MHG-GNN is 0.013 R² lower.
% - Keep robustness apart from accuracy. On QM9, NGBoost's highest noise robustness (normalised area under the R²-noise curve) is on Avalon, 0.982 against 0.975 for PDV.
% - Change: cut the sentence guessing that Avalon captures features relevant to efflux. Nothing tested it.
% - Change: typos. "overidden" should be "overridden"; "its own type" should be "their own type"; "data set" is "dataset" elsewhere.
% Paragraph 2: what a hybrid is, in words
% - An adaptive hybrid representation is a feature set built by selecting up to a fixed number of features (the budget) from several base representations of the same molecules.
% - It is chosen by a search, for one dataset and one downstream model. A different dataset, split or model gives a different hybrid.
% - It is an input, not a model. Any standard supervised model can be trained on it.
% - Its intended use: representations that say different things about the molecule (PDV, ECFP4, graph embeddings, sequence embeddings), with the features for this task picked from all of them.
% Paragraph 3: why select (her second paragraph, plus one sentence)
% - Pooling every feature adds irrelevant features that act as noise (Chapter~\ref{chap:flaws}).
% - Add one sentence: distance-based noise handling did worst in an earlier exploratory study (Section~\ref{sec:lnl_motivation}).
% - That study is motivation, not a result of this thesis. It covered one dataset (QM9 HOMO-LUMO gap) and one noise type (Gaussian).
% - The label-noise question is measured in Chapter~\ref{chap:similarity}, over more datasets and noise types.
% - Change: typos. "the feature which carry" should be "the features which carry"; "data set" is "dataset".
% Paragraph 4: why the budget is split per representation, and what the search has to decide
% - Representations differ in width by orders of magnitude, so wide ones can crowd out narrow ones in a single ranking.
% - Importance scores from different feature types are not on a common scale.
% - The search also includes one pooled allocation method, so per-representation and pooled selection are compared, not assumed.
% - One sentence listing what the idea leaves open:
%   - how useful each feature is; how many features each representation gets; how many representations take part;
%   - how a candidate hybrid is judged during the search; the total budget; whether the answer depends on the downstream model.
% Paragraph 5: two ways to use a hybrid for similarity
% - One sentence: both are defined in Chapter~\ref{chap:similarity}, "Similarity in a hybrid".
% - Figure: several base representation matrices of different widths; within each, features ranked by importance; the budget as a bar split between representations; the top features from each joined into a narrow hybrid. Side panel: the full concatenation, for contrast.
% ---- END NOTES ----

% TODO: is this right??? I'm making shit up with NGBoost.
Here, we present KIRBy (Knowledge Integration by Representation), a framework that takes a collection of feature matrices for a given data set and outputs a customised feature set optimised for a specific prediction task. The inputs can be any type of frozen feature vector, such as fingerprints, descriptor sets, foundation model embeddings, gene expression signatures, histopathology features. Within a given domain, these individual vectors typically contain its own type of information. It's a well-established phenomenon when it comes to machine learning, there's rarely ever a one-size-fits-all approach. Take the noise robustness study above as evidence, NGBoost's predictive performance on QM9 was only ever strong with given PDV input features. However, on Caco-2, NGBoost preferred Avalon, with PDV as a close second. The model NGBoost itself produces better fits when given physicochemical properties, yet this preference can be overidden by a different prediction task. Avalon likely captures features that are pertinent to the log efflux ratio.

% TODO: REVIEW MODIFICATION
Here, we present the adaptive hybrid representation, or ``hybrid'' for short. It takes a collection of feature matrices for a given dataset and outputs a customised feature set optimised for a specific prediction task. The inputs can be any type of frozen feature vector, such as fingerprints, descriptor sets, foundation model embeddings, gene expression signatures, histopathology features. Within a given domain, these individual vectors typically contain their own type of information. It's a well-established phenomenon when it comes to machine learning, there's rarely ever a one-size-fits-all approach. Take the noise robustness study in Chapter~\ref{chap:robustness} as evidence, with NGBoost on clean labels. On QM9, NGBoost's held-out $R^2$ was highest with PDV input features (0.865, median of 10 replicates). MHG-GNN was only 0.013 lower (0.852), and ECFP4 was lowest (0.706). On LogD, PDV was again highest (0.673), with MHG-GNN second (0.658). On hERG $K_i$, MHG-GNN was highest (0.511). PDV was second (0.456), the median of the 4 of 5 folds whose clean $R^2$ was at least 0.3. On Caco-2, NGBoost appeared to prefer Avalon (0.402), with ChemBERTa at 0.404. Both numbers are medians of only 2 of 5 folds; the other 3 fell below a clean $R^2$ of 0.3 and were excluded. Over all 5 folds, PDV is highest (0.370), ahead of Avalon (0.287) and ChemBERTa (0.256). The model NGBoost itself produces better fits when given physicochemical properties, yet this preference can be overridden by a different prediction task. Accuracy and robustness also pick different representations. On QM9 under Gaussian noise, NGBoost's highest noise robustness (\aucnorm, the normalised area under the $R^2$ against noise level curve) was on Avalon, 0.982 against 0.975 for PDV.

But as important as the selection of useful features is, it is also important to filter out the useless features. We could throw the kitchen sink at the models and pool together features across a wide array of representations. Not only would that pose computational challenges, but it also introduces a huge set of irrelevant features that function as noise in downstream applications (Chapter~\ref{chap:flaws}). With all of that in mind, I decided to develop a method that selects the best features for a given data set and model. In theory, it should effectively filter out noisy features and isolate the feature which carry the most relevant information.

% TODO: REVIEW MODIFICATION
But as important as the selection of useful features is, it is also important to filter out the useless features. We could throw the kitchen sink at the models and pool together features across a wide array of representations. Not only would that pose computational challenges, but it also introduces a huge set of irrelevant features that function as noise in downstream applications (Chapter~\ref{chap:flaws}). In an earlier exploratory study, the noise handling methods that used a distance between molecules did worst (Section~\ref{sec:lnl_motivation}). That study was motivation only: it covered one dataset, the QM9 HOMO--LUMO gap, and one noise type, Gaussian. Chapter~\ref{chap:similarity} measures the label-noise question over more datasets and noise types. With all of that in mind, I decided to develop a method that selects the best features for a given dataset and model. In theory, it should effectively filter out noisy features and isolate the features which carry the most relevant information.

% TODO: review, needs manual review
The questions this chapter and Chapter~\ref{chap:similarity} address are set out in Section~\ref{sec:intro:questions}.

% CODE INPUT
An adaptive hybrid representation is a feature set built by selecting up to a fixed number of features, the budget, from several base representations of the same molecules. It is chosen by a search, for one dataset and one downstream model. A different dataset, split or model gives a different hybrid. The hybrid is an input, not a model. Any standard supervised model can be trained on it. It is meant for base representations that say different things about the molecule, such as PDV descriptors, ECFP4 bits, graph embeddings and sequence embeddings. The features for the task are then picked from all of them.

% CODE INPUT
The budget is split between the representations before any feature is picked. Representations differ in width by orders of magnitude, from tens to thousands of features. In one ranking over all their features, a wide representation can crowd out a narrow one. Importance scores from different types of feature are also not on a common scale. So each representation's features are ranked only against each other. The search also includes one pooled allocation method, the global pool, which ranks the features offered by every representation together. Per-representation and pooled selection are therefore compared, not assumed. The idea leaves six choices open. Three concern the features: how useful each feature is, how many features each representation gets, and how many representations take part. The other three are how a candidate hybrid is judged during the search, the total budget, and whether the answer depends on the downstream model.

% CODE INPUT
A hybrid can also be used to measure similarity, in two ways. Distances can be computed in an ordinary hybrid, or the search can pick features by how well each molecule's activity is predicted from its five nearest neighbours. Both are defined in Chapter~\ref{chap:similarity}, ``Similarity in a hybrid''.


\subsection{Blueprints and what they tell a chemist}
% ---- NOTES: Blueprints and what they tell a chemist ----
% Takeaway: A finished search is saved as a blueprint, and every selected feature traces back to its representation, so the search says something about the problem, not only its score.
% Takeaway: Named features let a chemist look for new chemistry or confirm the selection makes sense. Learned embedding features cannot be read this way.
% Paragraph 1: the output is more than a score
% - A finished search is saved as a blueprint: the settings needed to rebuild the hybrid.
% - Every selected feature traces back to its representation and its importance rank.
% - So the search returns information about the problem, not only a better score.
% - The search keeps the best configuration for each downstream model and budget, so one dataset gives one blueprint per model and budget. Comparing them at one budget shows whether tree ensembles and neural networks draw on different features.
% Paragraph 2: which features a chemist can read
% - What a chemist can read: descriptor names (e.g. molecular weight, logP), the atom environment behind a selected ECFP bit, the substructure behind a MACCS key.
% - Outside small molecules the readable unit changes: a sequence position and residue for antibodies, a pathway for gene expression.
% - Learned embedding features say which representation was useful, not which chemistry. Chapter~\ref{chap:flaws} covers why.
% Paragraph 3: two directions, discovery and confirmation
% - Discovery: substructures or descriptors that recur among the kept features point to parts of the molecule that may be involved in binding or activity. A starting point for a chemist, not a finding.
% - Confirmation: if the kept features make chemical sense for the endpoint, that supports the claim that the selection finds real signal and is not fitting noise.
% - Caveat: selected means predictive, not causal. Correlated features can stand in for each other, so an unselected feature may matter just as much.
% - Caveat: a blueprint rebuilt on other molecules can select other features. Read features from the fit on the dataset in question.
% - Context for explainability in drug discovery: \citep{JimenezLuna2020}. Interpretable inputs as an alternative to explaining a black box after the fact: \citep{Rudin2019}.
% Paragraph 4: other uses of a task-specific similarity
% - A hybrid also gives a distance chosen for one task. Anything that uses nearby molecules as evidence could use it instead of a standard fingerprint distance.
% - One sentence naming the uses argued in Chapter~\ref{chap:flaws}: noise detection, nearest-neighbour prediction, applicability domain, activity-cliff detection, clustering molecules for group-aware splits, and weighting predictions from separate sites.
% - One use not covered there: the predictor that scores molecules in generative design \citep{Olivecrona2017}. This is future work.
% ---- END NOTES ----
\todo[inline, color=yellow!40]{5.1.4: explainability motivation; other applications of a task-specific similarity.}

% CODE INPUT
A finished search is saved as a blueprint: the settings needed to rebuild the hybrid. Every selected feature traces back to its representation and to its importance rank within that representation. So the search returns information about the problem, not only a better score. The search keeps the best configuration for each pair of downstream model and budget. One dataset therefore gives one blueprint per downstream model and budget. Comparing them at one budget shows whether tree ensembles and neural networks draw on different features.

% CODE INPUT
Some selected features can be read by a chemist. A descriptor from the PDV set carries its name, such as molecular weight or logP. A selected ECFP bit can be shown as the atom environment that sets it. A MACCS key stands for a known substructure. Outside small molecules the readable unit changes. For an antibody it is a sequence position and the residue at it. For gene expression it is a named gene set, such as a pathway or a chromosome arm. A feature of a learned embedding says only which representation was useful, not which chemistry. Chapter~\ref{chap:flaws} covers why.

% CODE INPUT
Readable features can be used in two directions. The first is discovery. Substructures or descriptors that recur among the kept features point to parts of the molecule that may be involved in binding or activity. They are a starting point for a chemist, not a finding. The second is confirmation. If the kept features make chemical sense for the endpoint, that supports the claim that the selection finds real signal and is not fitting noise. Explainability is a recognised need in drug discovery \citep{JimenezLuna2020}. Inputs that can be read directly are an alternative to explaining a black-box model after the fact \citep{Rudin2019}. Two caveats apply. First, selected means predictive, not causal. Correlated features can stand in for each other, so an unselected feature may matter just as much. Second, a blueprint rebuilt on other molecules can select other features. Features should be read from the fit on the dataset in question.

% CODE INPUT
A hybrid also gives a distance chosen for one task. Anything that uses nearby molecules as evidence could use it instead of a standard fingerprint distance. Chapter~\ref{chap:flaws} argues for several such uses: noise detection, nearest-neighbour prediction, applicability domain, activity-cliff detection, clustering molecules for group-aware splits, and weighting predictions from separate sites. One use is not covered there. Generative design scores each proposed molecule with a predictor \citep{Olivecrona2017}, and that predictor could take a hybrid as its input. We leave this to future work.


\subsection{Beyond small molecules}
% ---- NOTES: Beyond small molecules ----
% Takeaway: The search needs only a supervised task and several vector representations of the same samples, so it applies to tabular data in general.
% Takeaway: A dataset needs enough samples, and base representations whose features carry importance.
% Paragraph 1: what the method needs, and so where it applies
% - The search needs a supervised task and several vector representations of the same samples. Each representation is one feature matrix, one row per sample.
% - The default search objective scores prediction only. It reads no chemical structure.
% - So it applies to tabular data in general: any task where several encoders or measurements exist for the same inputs.
% - What does not carry over is chemistry-specific. Reading substructures from a blueprint needs named substructure features.
% - An optional search objective scores shared substructures between molecule pairs. It applies only to molecules.
% - The small-molecule distance and label-noise work of Chapter~\ref{chap:similarity} also does not carry over.
% Paragraph 2: one-hot encodings against learned sequence embeddings
% - One-hot sequence encodings are very sparse. At most positions the residue never changes across a dataset, so those features carry no importance.
% - One-hot can still work: it did on some antibody datasets.
% - Result needed: the antibody datasets where one-hot was among the selected representations, with the score in the dataset's metric.
% - Recommendation: prefer language-model sequence embeddings over raw one-hot encodings.
% Paragraph 3: the other domains used in this thesis, one or two sentences each
% - Antibodies: the most work after small molecules. Give this domain the most weight.
% - Chemical reactions (yield prediction): proof of concept. The representations are computed on reactants, products or their difference.
% - Nucleotide sequences (histone-mark and promoter tasks): proof of concept.
% - Cancer genomics (TCGA): each representation is a different measurement of the patient (expression, methylation, copy number, miRNA, protein, slide images), not a different encoder of one input. The search treats both cases the same way.
% Paragraph 4: what a dataset needs for the hybrid to help
% - Enough samples. The threshold (about 200 samples) waits on the full runs.
% - Base representations whose features carry importance. Where every representation already scores poorly, there are too few useful features to improve on.
% - Figure: one row per domain (small molecule, antibody, reaction, DNA sequence, patient sample). Left column: the raw input. Middle: several representation matrices. Right: the same search box, outputting a hybrid. Only the left column changes between rows.
% ---- END NOTES ----
\todo[inline, color=yellow!40]{5.1.5: other domains; requirements for a dataset.}

% CODE INPUT
The search needs two things: a supervised task, and several vector representations of the same samples. Each representation is one feature matrix, with one row per sample. The default search objective scores prediction only. It reads no chemical structure. So the method applies to tabular data in general: any task where several encoders or measurements exist for the same inputs. What does not carry over is specific to chemistry. Reading substructures from a blueprint needs named substructure features. An optional search objective scores how closely distances in the hybrid track the maximum common substructure shared by pairs of molecules. It applies only to molecules. The nearest-neighbour objective of Chapter~\ref{chap:similarity} needs only labels and a distance, so it does carry over. It is searched on the antibody and other non-molecular datasets as well.

% CODE INPUT
One-hot sequence encodings are very sparse. At most positions the residue never changes across a dataset, so those features carry no importance. One-hot can still work. Among the four antibody datasets of this chapter, one-hot is a base representation on nanobody thermostability only. The position-wise BLOSUM encoding is a base representation on humanness and trastuzumab--HER2 binding. One-hot was among the selected representations on \textbf{[RESULT NEEDED: the antibody datasets where one-hot was among the selected representations, with the score in the dataset's metric]}. We recommend language-model sequence embeddings over raw one-hot encodings.

% CODE INPUT
Antibodies received the most work after small molecules. This chapter uses four antibody datasets: Koenig expression, nanobody thermostability, humanness and trastuzumab--HER2 binding (Table~\ref{tab:datasets}). Their base representations are mostly embeddings from general protein language models and from antibody-specific language models. Some embeddings cover only the CDR loops, and some come from early or middle layers of ESM-2 (Table~\ref{tab:reps_antibodies}). Chemical reactions are a proof of concept, on two Buchwald--Hartwig yield datasets. Their base representations are computed on the reactants, on the products minus the reactants, or on the whole reaction; product fingerprints enter only as augmentation (Table~\ref{tab:reps_reactions}). Nucleotide sequences are also a proof of concept, on the H3K4me3 histone-mark and TATA promoter tasks (Table~\ref{tab:reps_nucleotides}). In cancer genomics (TCGA), each representation is a different measurement of the patient. These are gene expression, methylation (lung task only), copy number, miRNA, protein and slide images (Table~\ref{tab:reps_tcga}). In the other domains, each representation is a different encoder of one input. The search treats both cases the same way.

% CODE INPUT
A dataset needs two things for the hybrid to help. The first is enough samples. \textbf{[RESULT NEEDED: the number of samples below which the hybrid stops beating the best single representation; about 200 samples is expected]}. The second is base representations whose features carry importance. Where every representation already scores poorly, there are too few useful features to improve on.


